\documentclass{ieeeaccess}
\usepackage{cite}
\usepackage{amsmath,amssymb,amsfonts}
\usepackage{algorithmic}
\usepackage{graphicx}
\usepackage{textcomp}
\usepackage{booktabs}
\usepackage{multirow}
\usepackage{xcolor}
\usepackage{url}

% 中文支持（中文草稿阶段使用 XeLaTeX 编译）
\usepackage{xeCJK}
\setCJKmainfont{STSong}
\setCJKsansfont{STHeiti}
\setCJKmonofont{STFangsong}

\usepackage{bm}
\makeatletter
\AtBeginDocument{\DeclareMathVersion{bold}
\SetSymbolFont{operators}{bold}{T1}{times}{b}{n}
\SetSymbolFont{NewLetters}{bold}{T1}{times}{b}{it}
\SetMathAlphabet{\mathrm}{bold}{T1}{times}{b}{n}
\SetMathAlphabet{\mathit}{bold}{T1}{times}{b}{it}
\SetMathAlphabet{\mathbf}{bold}{T1}{times}{b}{n}
\SetMathAlphabet{\mathtt}{bold}{OT1}{pcr}{b}{n}
\SetSymbolFont{symbols}{bold}{OMS}{cmsy}{b}{n}
\renewcommand\boldmath{\@nomath\boldmath\mathversion{bold}}}
\makeatother

\def\BibTeX{{\rm B\kern-.05em{\sc i\kern-.025em b}\kern-.08em
    T\kern-.1667em\lower.7ex\hbox{E}\kern-.125emX}}

% 投稿占位符有意不标红，但仍需在视觉上明确显示。
% 生成终稿前请全文检索“占位”。
\newcommand{\TODO}[1]{\PackageWarning{manuscript}{Unresolved submission placeholder}\textbf{[占位：#1]}}

%% ============================================================
\begin{document}
\history{出版日期 xxxx 00, 0000，当前版本日期 xxxx 00, 0000。}
\doi{10.1109/ACCESS.2026.DOI}

\title{标称 3.91 ms UVC 设置下的时序像素掩蔽：一项传统 OCR 档位级测量研究}

\author{\uppercase{First A. Author}\authorrefmark{1}
\TODO{作者信息待填}}

\address[1]{\TODO{单位及地址待填}}

\tfootnote{\TODO{基金致谢待填。}
本文包含一项计划开展的最小风险人类被试研究。目前尚未
采集任何被试数据。只有在主管机构或院系出具
书面批准、豁免或免审认定后，方可开始招募；计划采用的保护措施
见第~\ref{sec:user_study} 节。}

\markboth
{作者 \headeretal：标称 UVC 设置下的时序像素掩蔽}
{作者 \headeretal：标称 UVC 设置下的时序像素掩蔽}

\corresp{通讯作者：\TODO{通讯作者邮箱}。}

%% ============================================================
%% 摘要
%% ============================================================
\begin{abstract}
屏幕内容可被静默拍摄，继而由光学字符识别（OCR）或视觉语言模型（VLM）处理。本研究测量在一个标称 USB 视频类（UVC）曝光控制设置下，传统 OCR 恢复率随档位发生的变化。该处理管线将每帧划分为高速交替的互补子帧，并加入随档位变化的叠加层。我们在 9 组距离--角度配置下共采集 10,575 张真实拍摄图像。主要的共同设置分析对三个档位分别使用同一组 288 个经重复采集平均的短曝光单元（12 个内容项 $\times$ 8 个拍摄几何条件 $\times$ 3 次重复）。按单次采集计算的引擎最优平均字符恢复率在未保护档、可读性优先档和高抑制档下分别为 94.5\%、17.8\% 和 5.6\%。在 432 个显式标注的字段评估机会中，可读性优先档仍有 17.7\% 的字段微平均完全恢复率，其匹配 P99 字符恢复率为 95.0\%。固定五变换、三引擎的网格最优选择（oracle）将可读性优先档和高抑制档的恢复率分别提高到 40.2\% 和 13.7\%。150 帧时域均值攻击和正轴 VLM 探针恢复出了多得多的文本。本实验没有确定有效曝光区间，也未将时序稀疏性与叠加层及占空比效应分离，因而不支持跨设备保护主张。本文的贡献是在单一 UVC 链路下获得的一项有明确边界的实证测量，以及足以排除部署主张或通用屏幕防拍摄主张的失效证据。
\end{abstract}

\begin{keywords}
对抗扰动，相机拍摄，显示隐私，光学字符识别，视觉暂留，屏幕--相机拍摄，时序掩蔽，视觉窃听
\end{keywords}

\titlepgskip=-21pt
\maketitle

%% ============================================================
%% I. 引言
%% ============================================================
\section{引言}

\PARstart{屏}{幕}遍布办公室、银行、医院和会议室；拍摄设备也同样随处可见，包括智能手机、智能眼镜~\cite{b_wang2026}（例如 Ray-Ban Meta）以及微型相机。仅凭一张屏幕照片，攻击者便可通过光学字符识别（OCR）提取文字，或通过目标检测识别界面元素。这种\emph{视觉窃听}十分隐蔽：它既不留下网络痕迹，也无需与目标设备发生物理接触。受控实验与实地研究均证实，此类攻击较为常见、成功率高，而且往往不会被察觉~\cite{b_ponemon2016,b_eiband2017}。尽管上述设备构成了本文的研究动机，但这里给出的物理实验证据仅覆盖一款 USB 视频类（UVC）网络摄像头；智能手机、智能眼镜及其他图像信号处理器（ISP）管线上的可复现性仍有待未来研究。

\subsection{现有方法及其局限}

现有屏幕隐私解决方案各自仅适用于有限场景。空间、软件、水印和时序方法分别针对不同的观察几何条件或攻击阶段。已有文献尚未报告与本文相同的场景组合，即正轴短曝光拍摄、静态敏感文本以及拍摄后的机器识别；详细比较见 \S\ref{sec:related}。因此，有必要针对机器从真实拍摄快照中恢复信息的能力开展专门测量。

难点在于，对相机端信息的削弱往往也会损害合法观看者的可读性。本文考察显示的时间结构能否在短曝光拍摄下扩大二者之间的差距。

\subsection{核心思想}

本文方法利用人类时间整合与相机曝光窗口之间的差异。每一帧被分解为 $n$ 个随机互补的子帧，并以高速交替显示。当曝光时间短于或接近一个子帧的持续时间时，相机只会记录不完整的像素子集，而人类观察者则会跨子帧整合内容（图~\ref{fig:concept}）。由于相机曝光并非瞬时采样，覆盖完整周期的长曝光或视频重构能够重新累加这些互补子帧；因此，本文不对任意曝光条件作不可恢复性保证。

本文将工作定位为受控测量研究，而非部署主张。我们考察在一个标称 3.91\,ms UVC 控制设置下，档位级显示处理如何改变传统 OCR 的恢复率。其操作场景是一条仪器化管线：每个目标由一张归档图像表示，随后再由传统 OCR 处理。该场景并非普通智能手机拍摄、智能眼镜拍摄，也不涉及能够针对性优化拍摄与识别过程的操作者。在 8 个共同设置的拍摄几何条件下，匹配平均恢复率在未保护档、可读性优先档和高抑制档下分别为 94.5\%、17.8\% 和 5.6\%。高抑制档未达到其描述性的 $<5\%$ 目标。可以成立的主张仅限于单一 UVC 条件下的这种档位级关联，而不是某个有效曝光区间、通用防拍摄设置，或时序稀疏性单独导致恢复率下降的证明。

边界证据本身也是测量贡献的一部分，而非附带说明。实验未记录环境照度，因而无法将日志中的 UVC 控制设置与典型智能手机自动曝光条件作强弱排序。所评估的时域均值攻击对 150 帧短曝光图像取平均，按 60\,fps 计算名义时长为 2.5\,s；本文未评估更短的视频持续时间。视觉语言模型（VLM）探针能恢复相当一部分大字号短片段，这表明对传统 OCR 的抑制无法迁移至现代端到端识别器（\S\ref{sec:vlm}）。

\begin{figure}[!t]
\centering
\includegraphics[width=\columnwidth,viewport=0 0 252 168,clip]{figures/visio/figure1_concept_threat/final/figure1_concept_threat.pdf}
\caption{时序像素掩蔽的威胁场景与运行边界。(a) 附近的相机无需访问主机系统，即可拍摄屏幕上的敏感内容。(b) 人类视觉会整合互补子帧，而相机的一次短曝光可能只与显示周期的一部分重叠，留给 OCR 的只是一个不完整片段。图中的相机路径仅为概念示意：它不假设零时长采样，也不假设固定的人类整合时间；长曝光或多帧重构能够恢复多得多的信息。}
\label{fig:concept}
\end{figure}

\subsection{贡献}

本文的主要贡献如下：
\begin{enumerate}
\item \textbf{受控的档位级物理测量}：在既有时序掩蔽思想的基础上，我们实现了像素级互补子帧、随档位变化的条纹/字形叠加层以及图形处理器（GPU）渲染原型。归档数据包含使用一款 eMeet SmartCam S600、在 3 个距离和 3 个角度下采集的 10,575 张图像。每个档位包含 288 个共同设置匹配单元（12 个内容项 $\times$ 8 个拍摄几何条件 $\times$ 3 次重复），其中未保护档、可读性优先档和高抑制档的平均恢复率分别为 94.5\%、17.8\% 和 5.6\%。这是在固定条件下，对经过几何校正的裁剪图像使用三个开源 OCR 引擎处理后得到的观察结果。

\item \textbf{考虑聚类的描述性证据与曝光敏感性}：我们报告采集层面的分布及按内容项聚类的配对对比。在归档数据的全部 9 个拍摄几何条件下，未保护短曝光相较于可读性优先档高出 78.0 个百分点（内容项聚类重采样区间 [75.5, 80.3]），相较于高抑制档高出 89.1 个百分点（[85.0, 92.4]）。将分析限制于 8 个共同设置的拍摄几何条件后，相应差值分别为 76.7 个百分点（[74.3, 79.2]）和 88.9 个百分点（[85.5, 91.9]）。这些区间描述结果对 12 个归档内容聚类的敏感性，并将所测试的拍摄几何条件视为固定；它们并不量化跨设备或跨场景的不确定性。

\item \textbf{失效边界刻画}：逐样本最强数字攻击的恢复率达到 95.2\%。对 150 帧相机图像取均值后，可读性优先档和高抑制档的恢复率分别达到 71.1\% 和 47.9\%；排除欠曝光的拍摄几何条件后，又分别升至 79.7\% 和 53.9\%。VLM 探针可恢复相当一部分大字号内容。这些负面结果界定了上述狭义传统 OCR 观察的终点：可读性优先档是唯一与计划开展的可用性协议相关联的档位，而高抑制档仍只是一个面向安全性的探索性压力测试档位，尚无用户可接受性证据。
\end{enumerate}

本文余下内容安排如下。\S\ref{sec:related} 回顾相关工作，\S\ref{sec:threat} 定义威胁模型，\S\ref{sec:method} 介绍系统设计。\S\ref{sec:evaluation} 报告主要 OCR 证据、VLM 与积分边界，以及延期执行的可用性协议。\S\S\ref{sec:discussion} 和~\ref{sec:conclusion} 综合证据层级并总结全文。检测/跟踪诊断被置于补充材料，因为它们并未扩展主要 OCR 主张。


%% ============================================================
%% II. 相关工作
%% ============================================================
\section{相关工作}
\label{sec:related}

\subsection{屏幕视觉窃听威胁}

已有研究记录了多种攻击，包括远距离屏幕读取攻击（望远镜、长焦镜头）~\cite{b_backes2009}、对 HDMI 电磁泄漏的深度学习重构~\cite{b_fernandez2024}、智能眼镜隐蔽拍摄~\cite{b_wang2026}，以及通过反射和折射实施的间接窃听~\cite{b_raguram2011}。Ponemon Institute 的一项受控研究表明，视觉入侵尝试的总体成功率为 91\%，其中 49\% 的情形在 15~分钟内出现首次成功入侵~\cite{b_ponemon2016}。Eiband 等~\cite{b_eiband2017} 调查了 174 名参与者，以研究肩窥的普遍程度及用户采用的应对策略。Tang 和 Shin~\cite{b_tang2023} 提出了 Eye-Shield：该方法利用前置摄像头检测旁观者，并实时模糊敏感区域，但它依赖用户侧硬件，且只能提供局部掩蔽。本文关注无需攻击者侧配合的显示侧缓解措施，但不声称能够在所有曝光和重构条件下阻止信息获取。

\subsection{显示隐私保护}

物理隐私滤镜能够限制可视角度，但无法阻挡正轴相机。Chen 等\ 提出 HideScreen，这是一种显示侧肩窥防御方案，通过抑制低频视觉线索，使屏幕内容在超出设计观看距离后与背景融为一体~\cite{b_hidescreen2019}；它针对空间域中的人类肩窥，而非短曝光相机拍摄或机器识别。Zhu 等~\cite{b_zhu2017} 提出了 LiShield，通过调制智能 LED 照明来干扰滚动快门相机，但该方案作用于环境照明，而非显示内容。视觉密码术将秘密图像拆分为多个份额，并在叠加时将其组合~\cite{b_naor1994}，由此启发了时序分解思想，但其原始形式面向印刷介质。数字水印~\cite{b_zhong2023} 可嵌入追踪标识，其中包括抗屏摄变体~\cite{b_fang2019}，但它提供的是事后归因，而非事前阻断。基于软件的数字版权管理（DRM）能够阻止数字截图，却无法防止物理拍摄。

屏幕--相机可见光通信（VLC）~\cite{b_nguyen2016,b_tran2021,b_ghiasi2024} 同样利用人类时间整合，将人眼不可见但可由相机解码的信息嵌入显示内容。本文方法与 VLC 互为镜像：VLC 旨在让相机\emph{读取}隐藏信息，而本文则旨在降低受限短曝光拍摄下显示内容对机器的可读性。

Yerazunis 和 Carbone 的早期工作~\cite{b_yerazunis2001} 探讨了用于显示隐私的时序图像掩蔽。Wu 和 Zhai~\cite{b_wu2013} 对时序心理视觉调制进行了系统化研究，Gao 等~\cite{b_gao2014} 则将其用于信息安全显示。Zhang 等~\cite{b_kaleido2015} 提出了 Kaleido，通过多帧分解对视频进行重新编码。Rainbow~\cite{b_rainbow2018} 则调制环境照明，以破坏滚动快门相机对实体物体的拍摄。这些方案使用的成像指标和主观质量指标，无法与本文的恢复率指标互换。既有研究没有报告本文所考察的特定组合：静态文本、短 UVC 曝光以及拍摄后的识别。Kaleido 使用色度互补帧与亮度变化，而本文的实现采用随机离散像素划分。探索性的高抑制档还放弃了严格完备性（\S\ref{sec:hardened}）。Lim~\cite{b_lim2007} 利用视觉暂留在浏览器中生成防屏幕截取键盘。表~\ref{tab:kaleido_compare} 将本研究置于有代表性的空间、照明和时序方案之中加以比较。

\begin{table*}[!t]
\caption{与代表性显示及实体场景隐私方案的适用范围比较。表中各方案针对不同威胁与评估目标，不构成参数匹配的性能主张。}
\label{tab:kaleido_compare}
\centering
\footnotesize
\setlength{\tabcolsep}{3.2pt}
\begin{tabular}{p{0.15\textwidth}p{0.15\textwidth}p{0.20\textwidth}p{0.25\textwidth}p{0.17\textwidth}}
\toprule
方案 & 保护介质 & 目标威胁 & 机制 & 已报告评估 \\
\midrule
HideScreen~\cite{b_hidescreen2019} & 移动显示屏 & 超出设计距离的人类肩窥 & 空间域低频线索抑制 & 人类观看距离/可读性 \\
LiShield~\cite{b_zhu2017} & 实体场景 & 滚动快门相机拍摄 & 调制环境 LED 照明 & 拍摄图像干扰 \\
Kaleido~\cite{b_kaleido2015} & 显示视频 & 连续录像盗录 & 感知完备的色度/亮度重编码 & 峰值信噪比（PSNR）、结构相似性指数（SSIM）和主观质量 \\
Rainbow~\cite{b_rainbow2018} & 实体物体 & 移动相机盗录 & 环境照明调制 & 滚动快门图像退化 \\
Lim~\cite{b_lim2007} & 浏览器键盘 & 单次数字屏幕截取 & 视觉暂留局部图案 & 键盘截取演示 \\
本文 & 静态屏幕文本 & 短曝光物理拍摄后的机器识别 & 随机像素划分；高抑制档破坏完备性 & OCR/VLM 恢复率与积分边界 \\
\bottomrule
\end{tabular}
\end{table*}

本文的贡献并非孤立的时序划分，而是在单一 UVC 条件下开展的系统级测量，并同时给出显式字段、积分及 VLM 失效证据。检测、跟踪和数字可用性代理仍属于补充诊断。

\subsection{对抗扰动与物理世界鲁棒性}

Wei 等~\cite{b_wei2024} 综述了数字对抗扰动~\cite{b_goodfellow2015,b_madry2018,b_carlini2017} 与物理世界对抗补丁~\cite{b_eykholt2018,b_thys2019,b_xu2020}。Athalye 等~\cite{b_athalye2018} 通过变换期望（Expectation over Transformation, EoT）首次实现了跨视角鲁棒的 3D 物理对抗物体，但该方法仍需对白盒目标模型具有访问权限。物理部署会引入数字评估中不存在的变换，因此必须在不同拍摄条件下测试鲁棒性，而不能从像素空间性能直接推断~\cite{b_eykholt2018}。物理对抗补丁具有跨视角鲁棒性，但需要针对特定分类器预先生成扰动图案~\cite{b_thys2019,b_xu2020}，对未见模型的迁移能力有限~\cite{b_gu2024}。屏幕--相机链路还会引入摩尔纹（moir\'{e}）伪影；Niu 等~\cite{b_niu2021} 已证明，这类伪影本身也可充当对抗扰动。基础机制在显示层对内容进行时序划分，而非优化模型特定的空间扰动。增强档位还会加入噪声及条纹/字形叠加层，其物理迁移效果由实验评估，而非预先假定。物理摄影本质上是一个采样过程，而防护机制则嵌入显示本身。


%% ============================================================
%% III. 威胁模型
%% ============================================================
\section{威胁模型}
\label{sec:threat}

\subsection{受保护资产与评估内容}

研究动机所涉及的资产包括密码、金融数据、个人身份信息（PII）、源代码和机密文件。实验并未表明所有此类资产都能得到保护。语料覆盖凭据、数字串、URL、代码、段落以及一个文档页面。我们在评分前对 12 个内容项中的 9 个标注了十八个高价值字段，并排除了普通叙述文本和偶然出现的章节编号。

\subsection{攻击者能力}

攻击能力并不构成严格的单调层级，而是由曝光、拍摄几何条件、时间采样和后处理能力共同决定。本文按以下范围报告结果：
\begin{itemize}
\item \textbf{主要范围}：在标称 3.91\,ms UVC 曝光控制设置下进行仪器化快照采集，随后使用 Surya、EasyOCR 和 Tesseract 逐帧识别。主要结果使用原始相机裁剪图像。另一个三引擎压力测试允许攻击者针对每次采集，在五种预先声明的确定性变换和原始输入中进行选择。该设置模拟的是固定 UVC 设置下的孤立图像和有界预处理网格，而非普通的自动曝光手机快照、覆盖完整显示周期的观测，或能够调整无约束增强管线的针对性操作者。实验未记录环境照度。这是一个狭窄的实测条件，并不代表能力最强的合理攻击者。
\item \textbf{实际失效边界}：长曝光、连续录像后的时域视频平均、离轴拍摄、已知周期下的相位搜索与通道重构、目标检测/关联任务，以及商业视觉语言模型的端到端识别（\S\ref{sec:vlm}）。所评估的视频攻击对 150 帧短曝光图像取平均，按 60\,fps 计算名义时长为 2.5\,s。它证明持续录像可以构成一种实际绕过方式，但没有证明可从任意四帧突发序列中恢复信息。本文量化这些风险，但仅将高抑制档结果解释为缓解效果，而非阻断保证。
\item \textbf{尚未得到充分覆盖}：其他相机/显示器组合、精确的滚动快门行对齐、真实场景中的时域超分辨率、以显示序列训练的学习式重构器，以及长期自适应攻击。
\end{itemize}

\subsection{攻击者知识}

攻击者\emph{知晓}时序掩蔽算法、子帧数~$n$ 以及候选周期参数，但\emph{不知道}当前周期的掩码种子或相位起点。密码学安全伪随机数生成器（CSPRNG）为每个周期提供掩码随机化，以防止固定模式学习和跨周期相关；但它不对全周期平均攻击提供密码学保证。一旦攻击者获取并配准一个完整周期，未知种子便无法阻止线性重构。稀疏采样可能破坏传统 OCR 所依赖的字形证据，但当前设计并未将二值化、分割或识别中的任何一个阶段确定为责任环节。端到端 VLM 从同一批归档片段中恢复出了大量文本（\S\ref{sec:vlm}），因此传统 OCR 观察结果不能迁移至这些模型。

\subsection{保护目标}

本文采用短曝光下的分层设计目标来解释所评估的档位。高抑制档以传统 OCR 字符恢复率 $<5\%$ 且完全匹配率 $=0\%$ 为目标；可读性优先档则以相对降低至少 80\% 且完全匹配率 $<1\%$ 为目标。这些论文层面的阈值并非预注册的安全保证。它们不适用于其他曝光条件、自适应预处理、VLM 或其他设备。目标检测在交并比为 0.50 时的平均精度均值（mAP50）变化属于补充诊断，而非通过标准。$\Delta E_{00}$ 和结构相似性（SSIM）数值是数字重构代理，并非真实面板感知或用户舒适度的证据。


%% ============================================================
%% IV. 系统设计
%% ============================================================
